%! AP Statistics — Unit 3, Lesson 3.3 — Group Activity
\documentclass[10pt]{article}
\usepackage{apstats-article}
\usepackage{apstats-boxes}

\begin{document}

\pageheader{Unit 3, Lesson 3.3}{Group Activity: Design Your Sample}
\namepartnerperiod

\vspace{0.1in}

\begin{scenariobox}[Context]{sky}
\small
A large urban school district has 8{,}000 students across 4 grade levels (9--12),
16 schools, and classes averaging 30 students each. The district wants to survey
students about their satisfaction with the school lunch program.
Your group will design a sampling plan using each method and evaluate which is best.
\end{scenariobox}

\vspace{0.1in}

\textbf{Part 1: Design a sampling plan using each method.}

\begin{practicebox}
\small
\textbf{A. Simple Random Sample (SRS) --- Target: 200 students}

\begin{enumerate}[leftmargin=1.5em, itemsep=6pt]
  \item How would you number the students? \blank{9cm}
  \item What device would you use to select 200? \blank{9cm}
  \item One advantage of SRS here: \blank{9cm}
  \item One disadvantage: \blank{9cm}
\end{enumerate}
\end{practicebox}

\vspace{0.1in}

\begin{practicebox}
\small
\textbf{B. Stratified Random Sample --- Target: 200 students}

\begin{enumerate}[leftmargin=1.5em, itemsep=6pt]
  \item What are the natural strata in this district? \blank{8cm}
  \item How many students would you sample from each stratum (be proportional)?\\
        Grade 9: \blank{2cm} \quad Grade 10: \blank{2cm} \quad Grade 11: \blank{2cm}
        \quad Grade 12: \blank{2cm}
  \item Why is stratifying by grade level useful here?\\
        \writeline\\[1pt]\writeline
\end{enumerate}
\end{practicebox}

\vspace{0.1in}

\begin{practicebox}
\small
\textbf{C. Cluster Sample --- Target: all students in selected schools}

\begin{enumerate}[leftmargin=1.5em, itemsep=6pt]
  \item The district has 16 schools. If you randomly select 4 schools, how many
        students would be in your sample (approximately)? \blank{4cm}
  \item Is a school a reasonable cluster? Explain.\\
        \writeline\\[1pt]\writeline
  \item One practical advantage of cluster sampling here: \blank{9cm}
\end{enumerate}
\end{practicebox}

\vspace{0.1in}

\begin{practicebox}
\small
\textbf{D. Systematic Random Sample --- Target: approximately 200 students}

\begin{enumerate}[leftmargin=1.5em, itemsep=6pt]
  \item The district has 8{,}000 students on a master list. To get 200, your sampling
        interval $k$ is: $k = $ \blank{2cm}
  \item How do you choose your random start? \blank{9cm}
  \item A potential problem with systematic sampling here: \blank{9cm}
\end{enumerate}
\end{practicebox}

\vspace{0.1in}

\textbf{Part 2: Compare and recommend.}

\begin{practicebox}
\small
\begin{enumerate}[leftmargin=1.5em, itemsep=8pt]
  \item Which method best ensures that each grade level is fairly represented?
        \blank{4cm} \quad Explain: \blank{7cm}
  \item Which method would be easiest to implement in a school setting?
        \blank{4cm} \quad Explain: \blank{7cm}
  \item Which method would your group recommend for this district? Justify in 2--3
        sentences.\\
        \writeline\\[1pt]\writeline\\[1pt]\writeline
\end{enumerate}
\end{practicebox}

\end{document}
